\documentclass[10pt]{article}

\usepackage[margin=0.82in]{geometry}
\usepackage{amsmath,amssymb,bm}
\usepackage{booktabs,array}
\usepackage[dvipsnames]{xcolor}
\usepackage{enumitem}
\usepackage{microtype}
\usepackage{needspace}
\usepackage{hyperref}

\definecolor{DeepBlue}{HTML}{17324D}
\definecolor{Accent}{HTML}{B54A35}
\definecolor{SoftGray}{HTML}{F3F5F7}
\hypersetup{colorlinks=true,linkcolor=DeepBlue,citecolor=Accent,urlcolor=DeepBlue}
\setlist{nosep,leftmargin=1.4em}
\setlength{\parindent}{0pt}
\setlength{\parskip}{0.45em}
\renewcommand{\arraystretch}{1.15}

\newcommand{\R}{\mathbb{R}}
\newcommand{\calE}{\mathcal{E}}
\newcommand{\calN}{\mathcal{N}}
\newcommand{\calL}{\mathcal{L}}
\newcommand{\vect}[1]{\bm{#1}}

\title{\vspace{-1.5em}{\LARGE\textbf{Transformers-ACE}}\par\vspace{0.35em}
{\large A Strictly Local Equivariant ACE-Attention Interatomic Potential}\par\vspace{0.4em}
{\normalsize Methods and Architecture Specification, Version 2}}
\author{Transformers-ACE project}
\date{June 2026}

\begin{document}
\maketitle

\begin{abstract}
Transformers-ACE is a local, energy-conserving machine-learned interatomic
potential that combines compressed atomic-cluster-expansion (ACE) density
correlations with equivariant neighbor-set attention. The architecture is built
for scalable molecular dynamics without atom-to-atom propagation of updated
hidden states. Periodic geometry enters through shifted displacement vectors,
chemical and angular edge densities are represented in irreducible
representations of $O(3)$, and learned Clebsch--Gordan products produce
correlations through maximum body order four. Attention queries depend on the
current center, whereas keys and equivariant values depend only on fixed local
edge features. A compact $C^2$ cutoff and a null channel in the attention
normalization make energy and forces continuous when an edge leaves the neighbor
list. A scalar total energy is the only predicted observable; forces and stress
are exact derivatives of that energy. This document specifies the equations
implemented by architecture version 2, the training objective, locality and
equivariance properties, checkpoint semantics, and validation requirements.
\end{abstract}

\section{Design scope}

The objective is an accurate short-range potential with linear system-size
scaling and conservative dynamics. The implementation follows four constraints:
\begin{enumerate}
  \item all geometric interactions are restricted to a finite cutoff;
  \item every intermediate feature has a defined $O(3)$ transformation law;
  \item no updated hidden state of atom $j$ is transmitted to atom $i$;
  \item energy is a scalar and all mechanical observables are its derivatives.
\end{enumerate}
The strictly local construction is motivated by local equivariant potentials
such as Allegro~\cite{musaelian2023}, while the representation algebra uses the
ACE and Euclidean-neural-network frameworks~\cite{drautz2019,geiger2022}.
Transformers-ACE is not an EquiformerV2 implementation: it retains compact
$\ell_{\max}=2$ channels and a local edge-attention mechanism rather than a deep
high-degree message-passing transformer~\cite{liao2023}.

\section{Periodic local environments}

Let atom $i$ have Cartesian row vector $\vect r_i$, atomic number $Z_i$, and a
periodic cell matrix $\vect h$. ASE supplies an integer image shift
$\vect S_{ij}\in\mathbb{Z}^3$ for every directed edge $j\rightarrow i$. The
periodic displacement and distance are
\begin{align}
  \vect r_{ij} &= \vect r_j-\vect r_i+\vect S_{ij}\vect h, \\
  d_{ij} &= \lVert\vect r_{ij}\rVert_2, \qquad
  \widehat{\vect r}_{ij}=\vect r_{ij}/d_{ij}.
\end{align}
The neighbor set is
\begin{equation}
  \calN(i)=\{j,\vect S_{ij}:d_{ij}<r_{\mathrm c}\}.
\end{equation}
The directed convention is important: spherical harmonics and forces are
evaluated from neighbor to center using the same shifted vector in training and
ASE inference.

\section{Smooth radial basis}

Architecture v2 uses the compact quintic envelope
\begin{equation}
 f_{\mathrm c}(d)=
 \begin{cases}
  1-10x^3+15x^4-6x^5, & 0\leq x<1,\\
  0, & x\geq1,
 \end{cases}
 \qquad x=d/r_{\mathrm c}.
 \label{eq:cutoff}
\end{equation}
It satisfies
\begin{equation}
 f_{\mathrm c}(r_{\mathrm c})=f'_{\mathrm c}(r_{\mathrm c})
 =f''_{\mathrm c}(r_{\mathrm c})=0,
\end{equation}
which is required for smooth forces, stress, phonons, and cell dynamics at the
neighbor-list boundary. The default radial basis is
\begin{equation}
 B_n(d)=f_{\mathrm c}(d)\sqrt{\frac{2}{r_{\mathrm c}}}
 \frac{\sin(n\pi d/r_{\mathrm c})}{d},\qquad n=1,\ldots,N_r,
 \label{eq:bessel}
\end{equation}
with the analytic $d\rightarrow0$ limit. A Gaussian basis multiplied by the
same envelope is also available. For the CsPbI$_3$ configuration,
$r_{\mathrm c}=6$~\AA{} and $N_r=12$.

\section{Equivariant representations}

A feature transforming in irrep $(\ell,p)$ is denoted
$x^{(\ell,p)}_m$, where $m=-\ell,\ldots,\ell$ and $p\in\{+1,-1\}$
is parity. Under $g\in O(3)$,
\begin{equation}
 x^{(\ell,p)}\mapsto D^{(\ell,p)}(g)x^{(\ell,p)}.
\end{equation}
Spherical harmonics $Y_{\ell m}(\widehat{\vect r}_{ij})$ carry parity
$p=(-1)^\ell$. Clebsch--Gordan tensor products combine two irreps according to
the triangle and parity rules,
\begin{equation}
 [x\otimes_{\mathrm{CG}}y]^{(\ell,p)}_m
 =\sum_{\ell_1m_1,\ell_2m_2}
 W^{\ell_1\ell_2\ell}_{c_1c_2c}
 C^{\ell m}_{\ell_1m_1,\ell_2m_2}
 x^{(\ell_1,p_1)}_{c_1m_1}y^{(\ell_2,p_2)}_{c_2m_2},
 \quad p=p_1p_2.
 \label{eq:cg}
\end{equation}
The coefficients $C$ are fixed by symmetry and the multiplicity mixing weights
$W$ are learned. The implementation uses e3nn tensor products~\cite{geiger2022}.

\section{Chemical-angular edge density}

Atomic numbers are mapped to scalar embeddings
$\vect e_i=\mathrm{Emb}(Z_i)\in\R^{C_0}$. A radial network maps
$\vect B(d_{ij})$ to edge-dependent tensor-product weights. The edge density is
\begin{equation}
 \vect a_{ij}=
 \mathrm{TP}_{\rho}\!\left(
   \vect e_j,
   \{Y_{\ell m}(\widehat{\vect r}_{ij})\}_{\ell=0}^{\ell_{\max}};
   \mathrm{MLP}_{r}[\vect B(d_{ij})]
 \right).
 \label{eq:edge-density}
\end{equation}
Because Eq.~\eqref{eq:bessel} includes $f_{\mathrm c}$, every component of
$\vect a_{ij}$ vanishes smoothly at the cutoff. The permutation-invariant local
density at center $i$ is
\begin{equation}
 \vect A_i=\sum_{j\in\calN(i)}\vect a_{ij}.
 \label{eq:density}
\end{equation}
For the default v2 model, the compact correlation irreps are
\begin{equation}
 16\times0e\;\oplus\;8\times1o\;\oplus\;4\times2e,
 \label{eq:corr-irreps}
\end{equation}
with total component dimension $16+3(8)+5(4)=60$.

\section{ACE-style body-ordered correlations}

The density $\vect A_i$ is a two-body center--neighbor representation. Learned
recursive Clebsch--Gordan products construct correlations up to maximum body
order four:
\begin{align}
 \vect C_i^{(2)} &= \vect A_i,\\
 \vect C_i^{(3)} &= \mathrm{TP}_{3}(\vect C_i^{(2)},\vect A_i),\\
 \vect C_i^{(4)} &= \mathrm{TP}_{4}(\vect C_i^{(3)},\vect A_i).
 \label{eq:ace-recursion}
\end{align}
Each tensor product has independent learned multiplicity weights. This replaces
the architecture-v1 all-ones contraction, which made nominal copies of the same
irrep rank deficient. Because each input density is already a sum over the
neighbor set, Eq.~\eqref{eq:ace-recursion} is invariant to neighbor ordering and
implicitly contains all repeated-index and distinct-neighbor terms admitted by
the channel and angular truncations.

The node irreps are
\begin{equation}
 64\times0e\;\oplus\;32\times1o\;\oplus\;16\times2e,
 \label{eq:node-irreps}
\end{equation}
with component dimension $240$. The initial center state is
\begin{equation}
 \vect h_i^{(0)}=
 \mathrm{Center}(\vect e_i)
 +\sum_{\nu=2}^{4}\mathrm{Linear}_{\nu}(\vect C_i^{(\nu)}).
 \label{eq:center-state}
\end{equation}
Only scalar slots receive the explicit center embedding. Thus $Z_i$ remains
identifiable even if two centers have identical neighbor densities.

The construction is a compressed, learned ACE density-correlation basis. It is
not claimed to be a formally complete, untruncated linear ACE basis. Increasing
$\ell_{\max}$, radial resolution, multiplicities, and correlation order expands
the representation systematically but changes computational cost.

\section{Strictly local equivariant attention}

The attention block is local to one center and its fixed edge features. It does
not read $\vect h_j^{(t)}$ from a sender atom. Scalar queries are generated from
the pre-normalized center state, while keys use only scalar components of the
fixed edge density:
\begin{align}
 \vect q_{ip} &= W^Q_p\,\mathrm{LN}(\vect h^{(0)}_{i,\ell=0}),\\
 \vect k_{ijp} &= W^K_p\,\mathrm{LN}(\vect a_{ij,\ell=0}).
\end{align}
For head $p$, the invariant logit is
\begin{equation}
 s_{ijp}=\frac{\vect q_{ip}^{\mathsf T}\vect k_{ijp}}{\sqrt{d_k}}
 +b_p(\vect B(d_{ij}))
 -\mathrm{softplus}(\lambda_p)d_{ij},
 \qquad \widetilde s_{ijp}=s_{ijp}/T.
 \label{eq:logit}
\end{equation}
All terms are rotational invariants. The distance term is learned independently
for each head and $T$ is the training-time attention temperature.

\subsection{Cutoff-null normalization}

Ordinary softmax can cancel a cutoff when the final neighbor approaches
$r_{\mathrm c}$, because numerator and denominator vanish together. V2 adds a
unit null/self channel:
\begin{equation}
 \alpha_{ijp}=
 \frac{f_{\mathrm c}(d_{ij})\exp(\widetilde s_{ijp})}
 {1+\sum_{k\in\calN(i)}f_{\mathrm c}(d_{ik})
 \exp(\widetilde s_{ikp})}.
 \label{eq:cutoff-attention}
\end{equation}
The code evaluates Eq.~\eqref{eq:cutoff-attention} with a per-center maximum
subtraction for numerical stability. The null channel guarantees
$\alpha_{ijp}\rightarrow0$ when an edge exits the neighbor list.

Equivariant values are fixed-edge linear maps,
\begin{equation}
 \vect v_{ijp}=\mathrm{Linear}^{V}_p(\vect a_{ij}),
\end{equation}
and the residual update is
\begin{equation}
 \vect u_i=\mathrm{Linear}^{O}\!\left[
 \frac{1}{H}\sum_{p=1}^{H}\sum_{j\in\calN(i)}
 \alpha_{ijp}\vect v_{ijp}\right],
 \qquad
 \widetilde{\vect h}_i=\vect h_i^{(0)}+\vect\gamma_{\mathrm{attn}}\odot\vect u_i.
 \label{eq:attention-update}
\end{equation}
Both $\alpha_{ijp}$ and $\vect a_{ij}$ contain the cutoff, so the attention
message vanishes at least quadratically in the envelope near $r_{\mathrm c}$.

\subsection{Angular feedback into scalar energy}

An equivariant linear map cannot directly map $\ell>0$ values to a scalar.
V2 therefore forms atom-local invariant squared norms after the attention
residual:
\begin{equation}
 n_{ic\ell}=\sum_{m=-\ell}^{\ell}
 \left|\widetilde h^{(\ell)}_{icm}\right|^2,
 \qquad \ell>0.
 \label{eq:norm}
\end{equation}
These invariants are concatenated with the scalar channels and passed through a
scalar feed-forward network:
\begin{equation}
 \vect h_{i,0}=\widetilde{\vect h}_{i,0}
 +\vect\gamma_{\mathrm{ffn}}\odot
 \mathrm{MLP}_{\mathrm{ffn}}
 \!\left(\mathrm{LN}[\widetilde{\vect h}_{i,0},\vect n_i]\right).
 \label{eq:ffn}
\end{equation}
Equation~\eqref{eq:norm} lets interference among directional neighbor values
affect energy while preserving $O(3)$ invariance. The default CsPbI$_3$ model
uses one attention block, two heads, layer-scale initialization $10^{-2}$, and
dropout probability $0.03$.

\section{Locality, equivariance, and scaling}

For each center $i$, Eqs.~\eqref{eq:edge-density}--\eqref{eq:ffn} depend only on
$Z_i$ and the raw species and geometry of edges in $\calN(i)$. A change to an
updated hidden state $\vect h_j$ cannot alter center $i$. Thus stacking the
implemented attention blocks does not increase the receptive field beyond
$r_{\mathrm c}$. The model is strictly local in the same operational sense as a
finite-cutoff many-body potential.

Spherical harmonics, equivariant linear maps, Clebsch--Gordan products, scalar
attention weights, and invariant norms preserve the transformation law of every
channel. The atomic energy is therefore invariant to global rotations,
reflections, translations, and atom indexing, while its gradient transforms as
a vector. Runtime and memory scale as
\begin{equation}
 \mathcal O(N\overline n\,C_{\mathrm{edge}})+\mathcal O(NC_{\mathrm{CG}}),
\end{equation}
where $\overline n$ is the average neighbor count. There is no iterative
center-to-center message passing.

\section{Energy, forces, and periodic stress}

Only scalar channels enter the atomic-energy readout,
\begin{equation}
 E_i=\mathrm{MLP}_{E}(\vect h_{i,0}),\qquad
 E=\sum_iE_i+E_{\mathrm{ref}}(\{Z_i\}).
 \label{eq:energy}
\end{equation}
The readout widths are $64\rightarrow64\rightarrow1$. The reference energy is
either a fitted per-species linear baseline or a constant mean energy per atom;
it is independent of positions and does not change forces or stress.

Forces are exact negative energy gradients,
\begin{equation}
 \vect F_i=-\frac{\partial E}{\partial\vect r_i}.
 \label{eq:force}
\end{equation}
No direct force head is used. Consequently the force field is conservative up
to floating-point and neighbor-list tolerances.

Periodic stress is obtained by differentiating with respect to six symmetric
strain parameters. Define
\begin{equation}
 \vect\epsilon(\vect\eta)=
 \begin{pmatrix}
 \eta_1&\eta_4&\eta_5\\
 \eta_4&\eta_2&\eta_6\\
 \eta_5&\eta_6&\eta_3
 \end{pmatrix},\qquad
 \vect F_{\epsilon}=\vect I+\vect\epsilon,
\end{equation}
and deform both positions and cell as row vectors,
\begin{equation}
 \vect r_i'=\vect r_i\vect F_{\epsilon},\qquad
 \vect h'=\vect h\vect F_{\epsilon}.
\end{equation}
At zero strain, normal and shear components are reconstructed as
\begin{align}
 \sigma_{aa}&=\frac{1}{V}\frac{\partial E}{\partial\eta_a},
 &&a=1,2,3,\\
 \sigma_{xy}&=\frac{1}{2V}\frac{\partial E}{\partial\eta_4},\quad
 \sigma_{xz}=\frac{1}{2V}\frac{\partial E}{\partial\eta_5},\quad
 \sigma_{yz}=\frac{1}{2V}\frac{\partial E}{\partial\eta_6}.
 \label{eq:stress}
\end{align}
The factor $1/2$ appears because one shear parameter changes two symmetric
matrix entries. $V=|\det\vect h'|$ is the deformed volume. This sign convention
matches ASE's $\sigma_{ab}=V^{-1}\partial E/\partial\epsilon_{ab}$ convention;
ASE cell filters supply the sign required for cell forces. The v2 architecture
retains the validated v1 stress target conversion, Voigt ordering
$(xx,yy,zz,yz,xz,xy)$, and strain-gradient implementation unchanged.

\section{Training objective}

For structure $s$ with $N_s$ atoms, the supervised losses are
\begin{align}
 \calL_E&=\frac{1}{|\mathcal B|}\sum_{s\in\mathcal B}
 \left[\frac{E_s-E_s^{\mathrm{DFT}}}{N_s}\right]^2,\\
 \calL_F&=\frac{1}{\sum_s3N_s}\sum_{s,i,a}
 \left(F_{sia}-F_{sia}^{\mathrm{DFT}}\right)^2,\\
 \calL_\sigma&=\frac{1}{6|\mathcal B_\sigma|}
 \sum_{s\in\mathcal B_\sigma}\sum_{v=1}^{6}
 \left(\sigma_{sv}-\sigma_{sv}^{\mathrm{DFT}}\right)^2,\\
 \calL&=w_E\calL_E+w_F\calL_F+w_\sigma(t)\calL_\sigma
 +w_{\mathrm{Sob}}\calL_{\mathrm{Sob}}.
 \label{eq:loss}
\end{align}
The light local linearization regularizer is
\begin{equation}
 \calL_{\mathrm{Sob}}=
 \left[E(\vect r+\vect\delta)-E(\vect r)
 +\sum_i\vect F_i\cdot\vect\delta_i\right]^2,
 \qquad \delta_{ia}\sim\mathcal N(0,\sigma_\delta^2),
\end{equation}
with detached forces in this auxiliary term. The default configuration uses
$w_E=1$, $w_F=10$, $w_\sigma:1000\rightarrow100000$ linearly over the first
20 epochs, $w_{\mathrm{Sob}}=10^{-3}$, and $\sigma_\delta=0.02$~\AA. The stress
ramp avoids introducing large second-derivative optimization terms at random
initialization.

Attention temperature decreases from $1.3$ to $1.0$ over 15 epochs. An optional
force-error adaptation multiplies this schedule by
$(\widehat{\calL}_F^{1/2}/F_{\mathrm{ref}})^\gamma$, with
$F_{\mathrm{ref}}=0.5$~eV/\AA{} and $\gamma=-0.25$ in the supplied CsPbI$_3$
configuration.

\Needspace{8\baselineskip}
\section{Data partitioning and optimization}

Ordered molecular-dynamics frames are strongly correlated. A frame-random split
can place adjacent configurations in training and validation, producing a
misleading generalization estimate. The default split selects complete
contiguous blocks of 25 frames for validation and removes a two-frame boundary
gap from training. Model initialization, block selection, and data-loader
shuffling use seed 42. An independent phase- and trajectory-grouped validation
file remains preferable when metadata are available.

The default optimizer is AMSGrad Adam with learning rate $10^{-3}$, weight decay
$10^{-5}$, gradient-norm clipping at 1, and cosine decay over 100 epochs to
$5\times10^{-5}$. Batch size is eight. Validation checkpointing begins after the
stress ramp; \texttt{model.pt} stores the best validation loss and
\texttt{model\_last.pt} stores the final state. Early stopping uses patience 20,
minimum epoch 25, and minimum improvement $10^{-4}$.

\section{Default CsPbI\texorpdfstring{$_3$}{3} architecture}

\begin{table}[h]
\centering
\small
\begin{tabular}{>{\raggedright\arraybackslash}p{0.44\linewidth}p{0.43\linewidth}}
\toprule
Quantity & Value \\
\midrule
Cutoff / radial basis & 6.0~\AA{} / 12 Bessel functions \\
Maximum angular degree & $\ell_{\max}=2$ \\
Node irreps & $64\times0e+32\times1o+16\times2e$ \\
Correlation irreps & $16\times0e+8\times1o+4\times2e$ \\
Maximum correlation/body order & 4 \\
Radial network & $12\rightarrow32\rightarrow$ TP weights \\
Attention & 1 local block, 2 heads \\
Scalar FFN & invariant scalars and $\ell>0$ squared norms \\
Readout & $64\rightarrow64\rightarrow1$ \\
Trainable parameters & 132,005 \\
Direct force/stress heads & none \\
Long-range electrostatics & not included \\
\bottomrule
\end{tabular}
\caption{Architecture-version-2 defaults used by the supplied CsPbI$_3$
training configuration.}
\end{table}

Relative to v1, the default v2 model reduces the parameter count from 593,097
to 132,005 by replacing the oversized radial/contraction parameterization with
compact independent channels and one local attention block.

\section{Verification and scientific use}

Automated regression tests cover:
\begin{itemize}
  \item sensitivity of a center representation to its own species;
  \item non-collapsed Clebsch--Gordan multiplicity channels;
  \item value, first derivative, and second derivative of the cutoff;
  \item energy and force continuity when an edge leaves the neighbor list;
  \item rotation, translation, and atom-permutation behavior;
  \item absence of updated sender-state dependence in attention;
  \item finite-difference agreement for all symmetric stress components;
  \item deterministic blocked validation splitting; and
  \item strict loading of both legacy-v1 and versioned-v2 checkpoints.
\end{itemize}

For scientific deployment, low test RMSE is necessary but insufficient. A
checkpoint should additionally reproduce phase-resolved errors, energy--volume
curves, elastic constants, phonons, defect geometries, and stable $NVE/NPT$
dynamics over the intended temperature and pressure range. Variable-cell
relaxations are meaningful only after the stress sign, units, and finite-strain
response have been validated against the reference electronic-structure method.

\section{Limitations}

The model is finite-cutoff and contains no explicit Coulomb, Ewald, charge
equilibration, dispersion-tail, or reciprocal-space term. This is intentional
for the present implementation. For ionic CsPbI$_3$, transferability with
respect to cutoff, cell size, dielectric environment, charged defects, and
long-wavelength distortions must therefore be tested explicitly. Long-range
physics should be added only when controlled cutoff and equation-of-state tests
demonstrate a deficiency that cannot be resolved by local data coverage.

The learned compressed correlations are not a proof of ACE basis completeness.
The current validation split reduces temporal leakage but cannot replace truly
independent trajectories and phases. Finally, architecture-v1 and v2 checkpoints
are not weight-compatible. Checkpoints without an
\texttt{architecture\_version} field are evaluated by the legacy model; new
training writes \texttt{architecture\_version: 2} and cannot resume v1 weights.

\section*{Code availability and reproducibility}

The implementation, training configuration, tests, and CsPbI$_3$ evaluation
workflow are maintained at
\url{https://github.com/paramvir3/Transformers-ACE}. The architecture equations
in this document correspond to the version-2 classes
\texttt{ACEV2Descriptor}, \texttt{StrictLocalEquivariantAttentionBlock}, and
\texttt{TransformersACE}. Exact package versions and checkpoint hashes should be
recorded for every reported result.

\begin{thebibliography}{9}
\bibitem{drautz2019}
R. Drautz,
``Atomic cluster expansion for accurate and transferable interatomic
potentials,'' \emph{Physical Review B} \textbf{99}, 014104 (2019).
\href{https://doi.org/10.1103/PhysRevB.99.014104}{doi:10.1103/PhysRevB.99.014104}.

\bibitem{geiger2022}
M. Geiger and T. Smidt,
``e3nn: Euclidean Neural Networks,'' arXiv:2207.09453 (2022).
\href{https://arxiv.org/abs/2207.09453}{arXiv:2207.09453}.

\bibitem{musaelian2023}
A. Musaelian, S. Batzner, A. Johansson, L. Sun, C. J. Owen,
M. Kornbluth, and B. Kozinsky,
``Learning local equivariant representations for large-scale atomistic
dynamics,'' \emph{Nature Communications} \textbf{14}, 579 (2023).
\href{https://doi.org/10.1038/s41467-023-36329-y}{doi:10.1038/s41467-023-36329-y}.

\bibitem{liao2023}
Y.-L. Liao, B. M. Wood, A. Das, and T. E. Smidt,
``EquiformerV2: Improved Equivariant Transformer for Scaling to Higher-Degree
Representations,'' arXiv:2306.12059 (2023).
\href{https://arxiv.org/abs/2306.12059}{arXiv:2306.12059}.
\end{thebibliography}

\end{document}
